% =====================================================================
% Anexo H: tablas y detalle de la justificacion de parametros (Seccion 3.6),
% movidos desde el capitulo 3 para mantener directo el cuerpo del documento.
% =====================================================================

\section*{Anexo H. Detalle de la justificación de los parámetros}
\label{sec:anexo-parametros}
\phantomsection\addcontentsline{toc}{section}{Anexo H. Detalle de la justificación de los parámetros}

\setcounter{figure}{0}\renewcommand{\thefigure}{H.\arabic{figure}}
\setcounter{table}{0}\renewcommand{\thetable}{H.\arabic{table}}

\subsection*{H.1. Función de cada paso y efecto de omitirlo}


Antes de discutir valores hay que justificar que cada paso exista. La
Tabla~\ref{tab:funcion} lo resume; el título de cada paso enlaza con su
descripción.

\begin{table}[!htbp]
\centering
\setstretch{1}\small
\setlength{\tabcolsep}{4pt}
\caption[Qué aporta cada paso y qué pasaría sin él]{Qué aporta cada paso y qué pasaría sin él. Elaboración propia.}
\label{tab:funcion}
\begin{tabularx}{\linewidth}{@{}>{\raggedright\arraybackslash\hsize=0.6\hsize}X>{\raggedright\arraybackslash\hsize=1.2\hsize}X>{\raggedright\arraybackslash\hsize=1.2\hsize}X@{}}
\toprule
\textbf{Paso} & \textbf{Qué aporta} & \textbf{Sin este paso} \\
\midrule
\multicolumn{3}{@{}l}{\textit{Unidad 1}}\\
\hyperref[sec:u1-p1]{1. Carga} & Limita la memoria y deja el color en la escala de Sonata & Sonata recibiría color hasta 65\,535 en vez de 255 \\
\hyperref[sec:u1-p2]{2. Centrado} & Conserva la precisión: en UTM, la precisión simple redondea $y$ a pasos de 1\,m; centrada, a micrómetros & Molduras y aristas se perderían por redondeo \\
\hyperref[sec:u1-p3]{3. Filtrado SOR} & Retira puntos sueltos antes de que se promedien en el vóxel o tuerzan una normal & Sonata recibiría posiciones y normales de superficies que no existen \\
\hyperref[sec:u1-p4]{4. Voxelización} & Iguala la densidad y reduce la nube para que quepa en la GPU & Las zonas cercanas al escáner dominarían y la nube no cabría \\
\hyperref[sec:u1-p5]{5. Normales} & Da a Sonata su tercera entrada, que el escáner no mide & Sonata no tendría su entrada completa \\
\hyperref[sec:u1-p6]{6. Escritura} & Separa las dos unidades y permite verificar repeticiones & Habría que preprocesar cada vez que corre Sonata \\
\midrule
\multicolumn{3}{@{}l}{\textit{Unidad 2}}\\
\hyperref[sec:u2-p2]{2. Modelo} & Aporta lo que Sonata aprendió sin etiquetas & Con pesos al azar, los descriptores no dirían nada de la escena \\
\hyperref[sec:u2-p3]{3. Preparación} & Pone la nube en la escala de Sonata y guarda el índice inverso & No se podrían devolver los descriptores a cada punto \\
\hyperref[sec:u2-p4]{4. Inferencia} & Produce $f_{\mathrm{ssl}}$, el descriptor que reemplaza a $f_{\mathrm{geo}}$ & No existe el aporte de la tesis \\
\hyperref[sec:u2-p5]{5. Escritura} & Ambas ramas se extraen una vez por sección y M4 lee sus archivos & Habría que repetir dos inferencias costosas y mezclar entornos incompatibles \\
\bottomrule
\end{tabularx}
\end{table}


\subsection*{H.2. Pruebas de los parámetros}


Cada prueba cambia un solo parámetro sobre la misma sección, deja fijos los
demás y mide la salida con cuatro indicadores:
\begin{enumerate}
\item \textbf{Rugosidad $r$}: distancia media de cada punto al plano de sus
vecinos. Baja si se limpia ruido, pero también si se borra detalle, por eso no
se lee sola.
\item \textbf{Dispersión angular}: cuánto difiere cada normal de las de sus
vecinos.
\item \textbf{ARI y NMI} (Ecs.~\eqref{eq:ari} y~\eqref{eq:nmi}): cuánto se
parecen los grupos de $k$-means de los descriptores a los de la configuración
de referencia. Cerca de 1, el parámetro no cambia lo que recibe el Módulo~4.
\item \textbf{Costo}: puntos que quedan, tiempo y memoria.
\end{enumerate}
Se elige el valor más barato a partir del cual los indicadores dejan de mejorar.
Si quitar un paso no cambia nada, se retira del procedimiento. La
Tabla~\ref{tab:pruebas} resume las seis pruebas.

\begin{table}[!htbp]
\centering
\setstretch{1}\small
\setlength{\tabcolsep}{4pt}
\caption[Pruebas de los parámetros]{Pruebas de los parámetros. Elaboración propia.}
\label{tab:pruebas}
\begin{tabularx}{\linewidth}{@{}l>{\raggedright\arraybackslash\hsize=0.9\hsize}X>{\raggedright\arraybackslash\hsize=0.9\hsize}X>{\raggedright\arraybackslash\hsize=1.2\hsize}X@{}}
\toprule
\textbf{Prueba} & \textbf{Qué cambia} & \textbf{Valores} & \textbf{Qué se espera ver} \\
\midrule
P0 & Tamaño de bloque & 1, 5 y 20 millones & Ninguna diferencia en $S$ \\
P1 & Filtrado: con o sin, $q$ y $\lambda$ & Sin filtrado; $q$ de 10 a 40; $\lambda$ de 1,5 a 3 & Si el filtrado mejora $r$ y la dispersión en las zonas ralas \\
P2 & Vóxel $\ell$ y rejilla $g$ & $\ell$ de 2, 3, 5 y 10\,cm; $g$ de 2 y 4\,cm & Si bajar de 5\,cm cambia los descriptores (ARI) y si la sección cabe en la GPU \\
P3 & Vecinos de las normales & 10, 20, 30 y 40 & Dispersión en paredes y en aristas; requiere separar este valor del $q$ del filtrado en el código \\
P4 & Precisión del \texttt{.npz} & Media y simple & Si la precisión media cambia los descriptores (ARI) \\
P5 & Tamaño de parche & 512, 1024 y 2048 & Si un parche menor ahorra memoria sin cambiar los descriptores \\
\bottomrule
\end{tabularx}
\end{table}


\subsection*{H.3. Replicación del procedimiento}


Las pruebas anteriores solo valen si cada configuración se puede repetir y dar
el mismo resultado. Esta subsección explica por qué el procedimiento de las
Unidades~1 y~2 es replicable, qué hay que hacer para replicarlo con otra
configuración y qué límites tiene esa replicación.

\paragraph*{Por qué es replicable.}
Hay cuatro razones, y cada una remite a la sección donde ya está descrita.
\begin{enumerate}
\item Ningún valor está escondido en el código. Cada parámetro de los dos
programas es una opción de la línea de órdenes con un valor por defecto
(\enlace{sec:u1-reproducibilidad}{Unidad~1} y Tabla~\ref{tab:replica}). La única
excepción es el tamaño de parche del codificador, fijo en el programa, que la
\enlace{sec:param-p5}{prueba~P5} convierte en opción.
\item La Unidad~1 no tiene pasos aleatorios
(\enlace{sec:u1-reproducibilidad}{verificación de la Unidad~1}); la
comprobación en dos máquinas está en la Sección~\ref{sec:res-u1-repro}.
\item Cada corrida deja constancia de lo que hizo: los metadatos de la
\enlace{sec:unidad1}{Unidad~1} (Paso~6) y los campos del archivo de
descriptores de la \enlace{sec:rama-ssl}{Unidad~2} (Paso~5).
\item El código, la instalación y los pesos están en el
\enlace{sec:hardware-repo}{repositorio}, y las versiones de las bibliotecas en
el \hyperref[sec:anexo-software]{Anexo~A}.
\end{enumerate}

\paragraph*{La única fuente de azar.}
La rutina de preparación de Sonata elige al azar un punto cuando varios caen en
la misma celda de la rejilla. Por eso dos corridas sobre la misma matriz $S$
podían dar descriptores algo distintos en esas celdas. Para eliminar esa
variación, el programa \texttt{extraer\_features.py} fija ahora una semilla
aleatoria, con valor 0 por defecto, que se elige con la opción
\texttt{--semilla} y queda guardada en el archivo. Que dos inferencias en la GPU
con la misma semilla den exactamente los mismos valores se comprueba repitiendo
una extracción y comparando los dos archivos.

\paragraph*{Cómo se replica una configuración.}
Replicar una configuración consiste en cinco pasos, los mismos para cualquiera
de las pruebas.
\begin{enumerate}
\item Fijar los valores de la configuración, uno por parámetro, a partir de la
Tabla~\ref{tab:replica}.
\item Ejecutar \texttt{preprocesar.py} con esos valores, con la orden de la
\enlace{sec:u1-reproducibilidad}{verificación de la Unidad~1}. En el clúster se
usa \texttt{preprocesar.sbatch}, que acepta el vóxel, los vecinos y el umbral
como variables de entorno (\hyperref[sec:anexo-software]{Anexo~A}).
\item Comprobar los conteos del archivo de metadatos como indica la
\enlace{sec:u1-reproducibilidad}{verificación de la Unidad~1}. Si no coinciden,
la configuración no se reprodujo y no tiene sentido seguir.
\item Ejecutar \texttt{extraer\_features.py} sobre esa misma matriz $S$, con la
rejilla, la precisión y la semilla de la configuración, con la orden del
\hyperref[sec:anexo-software]{Anexo~A}.
\item Comparar los descriptores con los de la condición de referencia, con el
ARI y el NMI del \enlace{sec:param-criterio}{criterio común}, y revisar las dos
salidas en el \enlace{sec:visor}{visor}, que además comprueba que descriptores y
puntos correspondan.
\end{enumerate}

\begin{table}[!htbp]
\centering
\setstretch{1}\small
\caption[Cómo cambiar cada parte de la configuración]{Cómo cambiar cada parte de la configuración y qué comprobar después. Elaboración propia.}
\label{tab:replica}
\begin{tabularx}{\linewidth}{@{}>{\raggedright\arraybackslash\hsize=0.8\hsize}X>{\raggedright\arraybackslash\hsize=1.1\hsize}X>{\raggedright\arraybackslash\hsize=1.1\hsize}X@{}}
\toprule
\textbf{Qué cambia} & \textbf{Cómo se cambia} & \textbf{Qué se comprueba} \\
\midrule
Otro escaneo & \texttt{--input} con su archivo LAS & Conteos de su archivo de metadatos \\
Tamaño de vóxel & \texttt{--voxel}, o \texttt{VOXEL} en el guion & \texttt{n\_voxel} y memoria en la GPU \\
Vecinos del filtrado & \texttt{--q}, o \texttt{Q} en el guion & \texttt{n\_sor} \\
Umbral del filtrado & \texttt{--lam}, o \texttt{LAM} en el guion & \texttt{n\_sor} \\
Rejilla de Sonata & \texttt{--grid} & Puntos retenidos por la rejilla \\
Precisión de los descriptores & \texttt{--fp32} & Tamaño del archivo, ARI y NMI \\
Azar de la rejilla & \texttt{--semilla} & Descriptores idénticos con la misma semilla \\
Otra máquina u otra versión de Open3D & Ninguna opción, mismos valores & Mismos conteos, aunque las filas salgan en otro orden \\
\bottomrule
\end{tabularx}
\end{table}

\paragraph*{Límites de la replicación.}
Hay tres. El primero es el orden de las filas
(\enlace{sec:u1-reproducibilidad}{verificación de la Unidad~1}). El segundo es la memoria. Una
configuración que deje muchos más representantes, por ejemplo un vóxel menor,
puede no caber en los 48\,GB de la GPU, y en ese caso la falta de memoria se
reporta como resultado de la configuración y no se reduce la nube para que
quepa. El tercero es el tamaño de la sección: si la memoria crece con el número
de puntos, las secciones más grandes no cabrán en una sola pasada y su
replicación depende de la división en bloques de la
Sección~\ref{sec:observables} o de la \enlace{sec:param-p5}{prueba~P5}.

